% Chapter 10: Roadmap, Delivery Phases & Pedagogical Takeaways

\begin{frame}{Project Roadmap}{Implementation Status vs Architectural Roadmap}
    \begin{block}{TaxonVision Delivery Phases}
        An honest, transparent breakdown of the current implementation status versus future architectural milestones.
        \begin{itemize}
            \item \textbf{Phase 1 (Foundation):} Complete. Pixi toolchain, DVC/S3 streaming, pre-commit gates, license compliance.
            \item \textbf{Phase 2 (Model \& MLOps):} Complete. Frozen feature extraction, Class-Balanced Loss, Optuna, Conformal safety, k3s, GHCR, Jenkins.
            \item \textbf{Phase 3 (Continuous Intelligence):} Target. Full Milvus MMKG, Vision-RAG, Label Studio retraining trigger.
        \end{itemize}
    \end{block}

    \begin{center}
    \begin{tikzpicture}[
        node distance=0.4cm and 0.4cm,
        box/.style={rectangle, draw, rounded corners, align=center, minimum height=1.1cm, minimum width=3.8cm, font=\scriptsize},
        arrow/.style={-stealth, thick},
        done/.style={fill=plantgreen!10, draw=plantgreen, thick},
        target/.style={fill=deepblue!10, draw=deepblue, thick}
    ]
        % Nodes
        \node[box, done] (phase1) {
            \textbf{Phase 1: Foundation}\\[1mm]
            \tiny Pixi, DVC, Licensing, Gates\\
            \tiny \textcolor{plantgreen}{\textbf{[Complete]}}
        };

        \node[box, done, right=of phase1] (phase2) {
            \textbf{Phase 2: Core MLOps}\\[1mm]
            \tiny Frozen Features, Optuna, Conformal,\\
            \tiny k3s, GHCR, Jenkins\\
            \tiny \textcolor{plantgreen}{\textbf{[Complete]}}
        };

        \node[box, target, right=of phase2] (phase3) {
            \textbf{Phase 3: Intelligence}\\[1mm]
            \tiny Milvus MMKG, Vision-RAG,\\
            \tiny Label Studio Retraining\\
            \tiny \textcolor{deepblue}{\textbf{[Target]}}
        };

        % Arrows
        \draw[arrow] (phase1) -- (phase2);
        \draw[arrow] (phase2) -- (phase3);

    \end{tikzpicture}
    \end{center}

    \script{Let's review our progress transparently. Phase 1 and Phase 2 are complete and battle-tested in our repository today: our reproducible Pixi environment, S3 zero-copy streaming, Class-Balanced Loss runner, conformal safety bounds, Jenkins container builds, and Kubernetes manifests are live. Phase 3—connecting the Milvus vector database into our active learning webhook loop—is our ongoing target architecture roadmap.}
    \qa{How do we ensure that new feature development does not regress existing Phase 1 and 2 capabilities?}{Every commit is protected by our pre-commit gates, automated test fixtures in \texttt{tests/}, and CI pipeline validation including license and complexity audits.}
\end{frame}

\begin{frame}{Pedagogical Takeaways}{Classic MLOps Anti-Patterns \& Engineering Fixes}
    \begin{block}{Bootcamp vs. Enterprise}
        Core pedagogical lessons for transitioning into MLOps, contrasting common junior anti-patterns with enterprise architectural solutions.
    \end{block}

    \begin{center}
    \begin{tikzpicture}[
        node distance=0.15cm and 0.5cm,
        box/.style={rectangle, draw, rounded corners, align=center, minimum height=0.75cm, minimum width=6.2cm, font=\scriptsize},
        bad/.style={fill=herbired!10, draw=herbired, thick},
        good/.style={fill=plantgreen!10, draw=plantgreen, thick},
        arrow/.style={-stealth, thick, line width=0.8mm, draw=black!30}
    ]
        % Row 1
        \node[box, bad] (bad1) {\textbf{Monolithic Notebook Sprawl}\\Unversioned, stateful Jupyter notebooks};
        \node[box, good, right=1cm of bad1] (good1) {\textbf{Modular Python Packages}\\Strict complexity limits, typed functions};
        \draw[arrow] (bad1) -- (good1);

        % Row 2
        \node[box, bad, below=of bad1] (bad2) {\textbf{Local Dataset Hoarding}\\Downloading terabytes to local SSDs};
        \node[box, good, right=1cm of bad2] (good2) {\textbf{Zero-Copy S3 Streaming}\\DVC pointers, Apache Arrow Parquet};
        \draw[arrow] (bad2) -- (good2);

        % Row 3
        \node[box, bad, below=of bad2] (bad3) {\textbf{Uncalibrated Softmax}\\Trusting raw model probabilities blindly};
        \node[box, good, right=1cm of bad3] (good3) {\textbf{Split Conformal Prediction}\\Mathematically guaranteed safety bounds};
        \draw[arrow] (bad3) -- (good3);

        % Row 4
        \node[box, bad, below=of bad3] (bad4) {\textbf{Unbudgeted Cloud GPUs}\\Renting unmanaged A100s indefinitely};
        \node[box, good, right=1cm of bad4] (good4) {\textbf{Bare-Metal GPU Time-Slicing}\\Strict memory ceilings on RTX 5070 Ti};
        \draw[arrow] (bad4) -- (good4);

    \end{tikzpicture}
    \end{center}

    \script{For all of you in this course, if you take away four key lessons from this project, make it these four: First, stop writing monolithic notebooks—use modular packages with strict complexity limits. Second, don't download multi-terabyte datasets locally—stream them and cache embeddings. Third, never trust raw softmax probabilities—use conformal prediction sets for safety. And fourth, don't waste thousands of dollars on unmanaged cloud GPUs—budget your VRAM and use time-slicing.}
    \qa{What is the fastest way for a junior data scientist to adopt these best practices?}{Start by enforcing automated pre-commit hooks with Ruff and Mypy, tracking datasets with DVC, and writing unit tests for your data transformation functions before touching model code.}
\end{frame}

\begin{frame}{Conclusion}{Core Architectural Principles \& System Vision}
    \begin{block}{The Pillars of Production Systems Engineering}
        TaxonVision is built on three uncompromising principles that define enterprise MLOps.
        \begin{itemize}
            \item \textbf{Reproducibility by Default:} Using Pixi environments, DVC data lineage, and declarative k3s manifests ensures every run is perfectly replicable.
            \item \textbf{Mathematical Safety under Uncertainty:} Split Conformal Prediction guarantees error bounds, refusing confident misclassification.
            \item \textbf{Hardware-Conscious Sustainability:} Bare-metal GPU time-slicing, INT8 ONNX quantization, and zero-allocation inference maximize throughput per watt.
        \end{itemize}
    \end{block}

    \vspace{0.4cm}
    \statusdone

    \script{To conclude: MLOps is not just about using trendy tools like Kubernetes or Docker. It is a systems engineering mindset. It means building models that are reproducible down to the exact commit, mathematically honest about what they don't know, and engineered to run reliably on physical hardware within budget. TaxonVision proves that with the right architecture, we can tackle the most complex biodiversity challenges on Earth with mathematical confidence and operational excellence. Thank you!}
    \qa{Where can students inspect the source code and documentation?}{The entire repository, including Zensical documentation, LaTeX strategy reports, and CI/CD pipelines, is maintained under strict version control in the taxon-vision-mlops repository.}
\end{frame}
